% WHAT ESCAPE CHARACTER A FAULT NAMES.
%
% A name written into a fault stands behind \escapechar, which is not
% always a backslash and need not be a character at all:
%
%   \escapechar=`\\   ! You can't use `\prevdepth' after \advance.
%   \escapechar=`|    ! You can't use `|prevdepth' after |advance.
%   \escapechar=-1    ! You can't use `prevdepth' after advance.
%
% Both halves of that message obey it: the name the document wrote and the
% primitive the reference names for itself.
%
% An ACTIVE CHARACTER has no escape character in front of it -- `Use of ~
% doesn't match its definition' -- and the NULL control sequence, which has
% no name, is written as the two control words that make it, each behind
% its own escape character.
%
% An escape character of nought is a NUL, and a NUL is written the way the
% printer writes one: a line break where it is the \newlinechar -- which it
% is by default -- and `^^@' where it is not.
%
% HSTeX wrote a backslash whatever \escapechar said, and wrote the null
% control sequence as that backslash alone. trip sets \escapechar to `|' on
% line 342 and writes `\advance\prevdepth' on line 410; its `\csname
% \endcsname' on line 402 is used where a macro wants a parameter text.
\catcode`\{=1 \catcode`\}=2
\tracingonline=1
\catcode`~=13
\message{[A with the ordinary escape character]}
\advance\prevdepth by 1pt
\message{[B with the escape character changed]}
\escapechar=`|
\advance\prevdepth by 1pt
\message{[C with no escape character at all]}
\escapechar=-1
\advance\prevdepth by 1pt
\message{[D the null control sequence, both ways]}
\escapechar=`|
\expandafter\def\csname\endcsname#1.{}
\expandafter\csname\endcsname x,
\escapechar=`\\
\expandafter\csname\endcsname x,
\message{[E an active character, which has none]}
\def~#1.{}
~x,
\escapechar=`|
~x,
\message{[F an escape character of nought, which is a NUL]}
\escapechar=0
\def\q#1.{}
\q x,
\message{[G and the same where the NUL is not the newline character]}
\newlinechar=`Y
\q x,
\message{[done]}
\end
